%% ch07 --- A number for chaos: the Lyapunov exponent
%%
%% Stub, generated by tools/gen_stubs.py from tools/chapters.json.
%% The brief below is the contract for this chapter: what it must argue, what
%% it must buy the reader, and what it must NOT take from its neighbours.
%% Writing the chapter means deleting the stub block below (the one that opens
%% on the next \chapter line) and replacing it with the chapter text -- after
%% which this file is never regenerated.
%% If the brief turns out to be wrong once the code has been run, change it in
%% the manifest and record the contradiction in CLAUDE.md under *Resolved
%% questions*.

\chapter{A number for chaos: the Lyapunov exponent}
\label{ch:lyapunov}

\chapterstub{%
Argument: the average rate at which nearby starts separate is a single
number, the Lyapunov exponent \dash{} the average of the logarithm of the
stretching factor per step, \code{lambda = mean of log|f prime(x[n])|}
\dash{} positive for chaos, negative for a system that forgets its
disturbances. For the logistic map at r = 4 it is exactly the natural
logarithm of 2, which the reader computes and checks. Its reciprocal is the
Lyapunov time, and the prediction horizon follows: steps until an error of
size e0 reaches a tolerance T is about \code{log(T/e0)/lambda}. The payoff
is the cruel arithmetic of that formula: measuring the start a thousand
times more precisely buys only \code{log(1000)/lambda} more steps \dash{}
about ten at r = 4 \dash{} so better instruments and bigger computers buy
forecasts that are LONGER BY A CONSTANT, not proportionally longer. The
exponent plotted against r shows chaos coming and going, with negative dips
at the periodic windows. Payoff: the reader can compute a Lyapunov exponent
from a time series of the map, turn it into a horizon, and explain why
weather forecasts improve by days per decade rather than by weeks. Traps to
elicit: ten times better data gives ten times longer forecasts; a bigger
computer will eventually make any system predictable. Listings: the exponent
by averaging; the horizon table against initial error; the exponent across
r. Labs: compute the exponent for a given r; compute a prediction horizon.
Figures: lambda against r; horizon against the logarithm of initial error, a
straight line.

\medskip
\textbf{Word budget:} 3000.
}
